\subsection{MODULES OVER INTEGRAL DOMAINS}

PROPOSITION 23. In a module E over an integral domain A, the set T of non-free elements is a submodule of E.

If x and y are not free, there exist two non-zero elements \(\alpha\) , \(\beta\) in A such that \(\alpha x = 0\) and \(\beta y = 0\) . Then \(\alpha\beta \neq 0\) since A is an integral domain and \(\alpha\beta(\lambda x + \mu y) = 0\) for all \(\lambda\) and \(\mu\) in A since A is commutative, hence \(\lambda x + \mu y\) is not free.

Remark. Let E be a module over any commutative ring A. If x is a non-free element of E, every element of the submodule Ax is non-free. On the other hand, if A contains divisors of 0, the sum of two non-free elements of E may be free; for example, in Z/6Z considered as a module over itself, 3 and 4 are not free, but \(3 + 4 = 1\) is free.

Proposition 23 leads to the following definition:

DEFINITION 5. In a module E over an integral domain A, the torsion submodule of E is the submodule of E consisting of the non-free elements (also called the torsion elements of E).

When E is equal to its torsion submodule (that is when every element of E is annihilated by an element \(\neq0\) of A) E is called a torsion module. When the torsion submodule of E is reduced to 0 (that is every non-zero element of E is free) E is called (by an abuse of language) a torsion-free module.

Every submodule of a free A-module (and in particular every projective A-module) is torsion-free. The Z-module Q is torsion-free.

PROPOSITION 24. Let A be an integral domain. For every A-module E, let T(E) denote the torsion submodule of E. Let \(f: \mathbf{E} \to \mathbf{E}'\) be an A-linear mapping, E and \(\mathbf{E}'\) being A-modules.

\begin{enumerate}
\def\labelenumi{(\roman{enumi})}
\item
  \(f(\mathbf{T}(\mathbf{E})) \subset \mathbf{T}(\mathbf{E}')\) .
\item
  If \(f\) is injective, \(f(\mathbf{T}(\mathbf{E})) = \mathbf{T}(\mathbf{E}') \cap f(\mathbf{E})\) .
\item
  If \(f\) is surjective and \(\operatorname{Ker}(f) \subset \mathrm{T}(\mathbf{E})\) , then \(f(\mathrm{T}(\mathbf{E})) = \mathrm{T}(\mathbf{E}')\) .
\end{enumerate}

Assertions (i) and (ii) are obvious. On the other hand, if \(f\) is surjective and \(x' \in \mathrm{T}(\mathrm{E}')\) , then \(x' = f(x)\) , where \(x \in \mathbf{E}\) , and by hypothesis there exists \(\alpha \neq 0\) in \(\mathbf{A}\) such that \(f(\alpha x) = \alpha x' = 0\) ; whence \(\alpha x \in \operatorname{Ker}(f)\) and by the hypothesis there exists \(\beta \neq 0\) in \(\mathbf{A}\) such that \(\beta(\alpha x) = 0\) ; as \(\beta \alpha \neq 0\) , \(x \in \mathrm{T}(\mathrm{E})\) .

COROLLARY 1. For every A-module E, E/T(E) is torsion-free.

If \(f: \mathbf{E} \to \mathbf{E}'\) is an A-linear mapping, let \(f_{\mathrm{T}}\) denote the mapping \(\mathrm{T}(\mathbf{E}) \to \mathrm{T}(\mathbf{E}')\) with the same graph as the restriction of \(f\) to \(\mathrm{T}(\mathbf{E})\) . With this notation:

COROLLARY 2. For every exact sequence of A-modules and A-linear mappings

\[
0 \longrightarrow \mathrm{E} ^ {\prime} \stackrel {f} {\longrightarrow} \mathrm{E} \stackrel {g} {\longrightarrow} \mathrm{E} ^ {\prime \prime}
\]

the sequence

\[
0 \longrightarrow \mathrm{T} (\mathrm{E} ^ {\prime}) \xrightarrow {f _ {\mathrm{T}}} \mathrm{T} (\mathrm{E}) \xrightarrow {g _ {\mathrm{T}}} \mathrm{T} (\mathrm{E} ^ {\prime \prime})
\]

is exact.

\[
\operatorname{Ker} \left(g _ {\mathrm{T}}\right) = \operatorname{Ker} (g) \cap \mathrm{T} (\mathrm{E}) = f \left(\mathrm{E} ^ {\prime}\right) \cap \mathrm{T} (\mathrm{E}) = f \left(\mathrm{T} \left(\mathrm{E} ^ {\prime}\right)\right) = \operatorname{Im} \left(f _ {\mathrm{T}}\right).
\]

PROPOSITION 25. Let A be an integral domain and \((\mathbf{E}_{\iota})\) a family of A-modules; then

\begin{enumerate}
\def\labelenumi{(\arabic{enumi})}
\setcounter{enumi}{36}
\tightlist
\item
\end{enumerate}

\[
\mathrm{T} \left(\bigoplus_ {\iota} \mathrm{E} _ {\iota}\right) = \bigoplus_ {\iota} \mathrm{T} (\mathrm{E} _ {\iota}).
\]

Let \((x_{\iota})\) be an element of \(\bigoplus_{\iota} E_{\iota}\) such that \(x_{\iota} \in T(E_{\iota})\) for all \(\iota\) ; then each of the \(x_{\iota}\) is annihilated by an element \(\alpha_{\iota} \neq 0\) of A and it may be assumed that \(\alpha_{\iota} = 1\) when \(x_{\iota} = 0\) ; as the family \((x_{\iota})\) has finite support, the element \(\alpha = \prod_{\iota} \alpha_{\iota}\) of A is defined and \(\neq 0\) ; it obviously annihilates \(\bigoplus_{\iota} x_{\iota}\) and hence \(\bigoplus_{\iota} T(E_{\iota}) \subset T\left(\bigoplus_{\iota} E_{\iota}\right)\) ; the converse is immediate.

If E and F are two A-modules, clearly \(T(E \otimes_{A} F)\) contains the canonical images of \(T(E) \otimes_{A} F\) and \(E \otimes_{A} T(F)\) ; but examples can be given of torsion-free A-modules E, F such that \(T(E \otimes_{A} F) \neq 0\) (Exercise 31).

Note that an infinite product of torsion modules is not necessarily a torsion module; for example, in the Z-module \(\prod_{n=1}^{\infty}\left(\mathbf{Z}/p^{n}\mathbf{Z}\right)\) (p an integer \textgreater1), the element all of whose coordinates are 1 is free.

PROPOSITION 26. Let A be an integral domain, K its field of fractions, E an A-module and \(\mathbf{E}_{(\mathbf{K})} = \mathbf{K} \otimes_{\mathbf{A}} \mathbf{E}\) the vector space over K obtained by extending the ring of operators; let \(\phi\) denote the canonical A-linear mapping \(x \mapsto 1 \otimes x\) of E into \(\mathbf{E}_{(\mathbf{K})}\) .

\begin{enumerate}
\def\labelenumi{(\roman{enumi})}
\item
  Every element of \(\mathbf{E}_{(\mathbf{K})}\) is of the form \(\lambda^{-1}\phi(x)\) for \(\lambda \in \mathbf{A}\) , \(\lambda \neq 0\) and \(x \in \mathbf{E}\) .
\item
  The kernel of \(\phi\) is the torsion submodule T(E) of E.
\item
  Every element of \(\mathrm{E}_{(\mathbb{K})}\) is of the form \(z = \sum_{i=1}^{n} \xi_i \phi(x_i)\) with \(\xi_i \in \mathbb{K}\) and \(x_{i}\in \mathbf{E};\) for all \(i\) , there exists \(\alpha_{i}\in \mathbf{A}\) such that \(\alpha_{i}\neq 0\) and \(\alpha_{i}\xi_{i}\in \mathbf{A}\) ; if \(\alpha = \prod_{i = 1}^{n}\alpha_{i},\) then \(\alpha \neq 0\) and \(\alpha \xi_{i} = \beta_{i}\in \mathbf{A}\) for all \(i\) , whence, in \(\mathrm{E}_{(\mathbf{K})}\) ,
\end{enumerate}

\[
z = \alpha^ {- 1} (\alpha z) = \alpha^ {- 1} \sum_ {i = 1} ^ {n} \beta_ {i} \phi (x _ {i}) = \alpha^ {- 1} \phi \left(\sum_ {i = 1} ^ {n} \beta_ {i} x _ {i}\right)
\]

since \(\phi\) is A-linear.

\begin{enumerate}
\def\labelenumi{(\roman{enumi})}
\setcounter{enumi}{1}
\tightlist
\item
  If \(x \neq 0\) is not free in \(\mathbf{E}\) , there exists \(\alpha \neq 0\) in \(\mathbf{A}\) such that \(\alpha x = 0\) , whence \(\alpha \phi(x) = \phi(\alpha x) = 0\) in \(\mathrm{E}_{(\mathbb{K})}\) , which implies \(\phi(x) = 0\) . Conversely, suppose that, for some \(x \in \mathbf{E}\) , \(1 \otimes x = 0\) in \(\mathrm{E}_{(\mathbb{K})}\) ; we show that \(x\) is a torsion element of \(\mathbf{E}\) . We consider the set \(\mathfrak{M}\) of monogenous sub-A-modules of \(\mathbf{K}\) ; this is a right directed set under the relation of inclusion, for any two elements \(\alpha, \beta\) of \(\mathbf{K}\) can be written as \(\alpha = \zeta^{-1}\xi, \beta = \zeta^{-1}\eta\) , where \(\xi, \eta, \zeta\) belong to \(\mathbf{A}\) and \(\zeta \neq 0\) , hence \(\mathbf{A}. \alpha \subset \mathbf{A}. \zeta^{-1}\) and \(\mathbf{A}. \beta \subset \mathbf{A}. \zeta^{-1}\) . Moreover \(\mathbf{K}\) is the union of the modules \(\mathbf{M} \in \mathfrak{M}\) and can therefore be considered as the direct limit of the direct system defined by the modules \(\mathbf{M} \in \mathfrak{M}\) and the canonical injections (§ 6, no. 2, Remark). Hence also, to within a canonical isomorphism, \(\mathrm{E}_{(\mathbb{K})} = \varinjlim (\mathbf{M} \otimes_{\mathbf{A}} \mathbf{E})\) (§ 6, no. 3, Proposition 7) and the relation \(1 \otimes x = 0\) in \(\mathrm{E}_{(\mathbb{K})}\) implies that there exists an
\end{enumerate}

\(M \in \mathfrak{M}\) such that \(1 \in M\) and \(1 \otimes x = 0\) in the tensor product \(M \otimes_{A} E\) (Set Theory, III, § 7, no. 5, Lemma 1). It may further be supposed (replacing if need be \(M\) by a monogenous submodule \(M' \supset M\) of \(K\) ) that \(M = A. \gamma^{-1}\) , where \(\gamma \in A\) and \(\gamma \neq 0\) . Now the mapping \(\xi \mapsto \gamma\xi\) is an isomorphism of \(M\) onto the A-module \(A\) ; on the other hand, the canonical isomorphism \(A \otimes_{A} E \to E\) (§ 3, no. 4, Proposition 4) maps \(\xi \otimes x\) to the element \(\xi x\) of \(E\) ; thus there exists an isomorphism \(M \otimes_{A} E \to E\) which maps the tensor product \(\xi \otimes x\) to the element \((\gamma\xi)x\) of \(E\) . The hypothesis \(1 \otimes x = 0\) in \(M \otimes_{A} E\) thus implies \(\gamma x = 0\) .

Remark. Let \(\alpha^{-1}\phi(x)\) , \(\beta^{-1}\phi(y)\) be two elements of \(\mathbf{E}_{(\mathbb{K})}\) , with \(\alpha\in\mathbf{A}\) , \(\beta\in\mathbf{A}\) , \(x\in\mathbf{E}\) , \(y\in\mathbf{E}\) , \(\alpha\beta\neq0\) . In order that \(\alpha^{-1}\phi(x)=\beta^{-1}\phi(y)\) , it is necessary and sufficient that \(\beta x-\alpha y\) be a torsion element of \(\mathbf{E}\) , for this relation is equivalent to \(\beta\phi(x)=\alpha\phi(y)\) , which may also be written \(\phi(\beta x-\alpha y)=0\) .

COROLLARY 1. If E is a torsion-free A-module, the canonical mapping \(\phi: \mathbf{E} \to \mathbf{E}_{(\mathbb{K})}\) is injective.

Recall (§ 5, no. 1) that for every A-linear mapping of E into a vector space F over K, there exists one and only one K-linear mapping \(\tilde{f}\colon\mathbf{E}_{(\mathbb{K})}\to\mathbf{F}\) such that \(f=\tilde{f}\circ\phi\) ; we shall say that \(\tilde{f}\) is associated with \(f\) .

COROLLARY 2. Let \(f\) be an A-linear mapping of \(\mathbf{E}\) into a vector space \(\mathbf{F}\) over \(\mathbf{K}\) ; if \(\operatorname{Ker}(f) \subset \mathrm{T}(\mathbf{E})\) , the K-linear mapping \(\bar{f}\) associated with \(f\) is injective.

We write an element of \(\operatorname{Ker}(\bar{f})\) in the form \(\lambda^{-1}\phi(x)\) , where \(\lambda \in \mathbf{A}\) , \(\lambda \neq 0\) , \(x \in \mathbf{E}\) ; the relation \(\bar{f}(\lambda^{-1}\phi(x)) = 0\) is equivalent to \(\lambda^{-1}\bar{f}(\phi(x)) = 0\) in \(\mathbf{F}\) and hence to \(f(x) = \bar{f}(\phi(x)) = 0\) . By hypothesis, this implies \(x \in \mathrm{T}(\mathbf{E})\) and hence \(\phi(x) = 0\) , which proves the corollary.

COROLLARY 3. Let E be an A-module and g an A-linear mapping of E into a vector space F over K such that \(g(\mathbf{E})\) generates F and \(\operatorname{Ker}(g) \subset \mathrm{T}(\mathbf{E})\) . Then the K-linear mapping \(\bar{g}\) associated with g is an isomorphism of \(\mathrm{E}_{(\mathbb{K})}\) onto F.

\(\bar{g}\) is injective by Corollary 2 and the hypothesis that \(g(\mathbf{E})\) generates F implies that \(\bar{g}\) is surjective.

For every A-module E, the vector space \(E_{(K)}\) is said to be associated with E. For every subset S of E the rank of S over K (or by an abuse of language, the rank of S) is the rank of the canonical image \(\phi(S)\) of S in \(E_{(K)}\) , in other words (no. 2, Definition 1) the dimension over K of the vector subspace of \(E_{(K)}\) generated by \(\phi(S)\) .

When E is a torsion-free A-module, it is usually identified with its canonical image \(\phi(\mathbf{E})\) in \(\mathrm{E}_{(\mathbb{K})}\) . With this convention, every generating system of E contains a basis of \(\mathrm{E}_{(\mathbb{K})}\) (no. 1, Theorem 2). In particular:

COROLLARY 4. Every finitely generated A-module is of finite rank.

Note that the converse of this corollary is not necessarily true; for example Q is a Z-module of rank 1 but is not finitely generated over Z.

Recall (§ 5, no. 1) that for every linear mapping \(f: \mathbf{E} \to \mathbf{E}'\) (where \(\mathbf{E}\) and \(\mathbf{E}'\) are A-modules), \(f_{(\mathbf{K})}\) denotes the K-linear mapping \(1_{\mathbf{K}} \otimes f: \mathbf{E}_{(\mathbf{K})} \to \mathbf{E}_{(\mathbf{K})}'\) .

PROPOSITION 27. For every exact sequence

\[
\mathrm{E} ^ {\prime} \xrightarrow {f} \mathrm{E} \xrightarrow {g} \mathrm{E} ^ {\prime \prime}
\]

of A-linear mappings, the corresponding sequence of K-linear mappings

\[
\mathbf {E} _ {(\mathbf {K})} ^ {\prime} \xrightarrow {f _ {(\mathbf {K})}} \mathbf {E} _ {(\mathbf {K})} \xrightarrow {g _ {(\mathbf {K})}} \mathbf {E} _ {(\mathbf {K})} ^ {\prime \prime}
\]

is exact.

Suppose that \(g_{(\mathbb{K})}(\lambda^{-1} \otimes x) = 0\) , with \(\lambda \in \mathbf{A}\) , \(\lambda \neq 0\) , \(x \in \mathbb{E}\) ; this is equivalent to \(\lambda^{-1} \otimes g(x) = 0\) in \(\mathrm{E}_{(\mathbb{K})}'\) and hence also to

\[
1 \otimes g (x) = \lambda (\lambda^ {- 1} \otimes g (x)) = 0;
\]

by Proposition 26, there exists \(\alpha \neq 0\) in A such that \(\alpha g(x) = 0\) in \(\mathbf{E}^{\prime \prime}\) , or also \(g(\alpha x) = 0\) . By hypothesis, there is therefore an \(x' \in \mathbf{E}'\) such that \(\alpha x = f(x')\) and therefore \(\lambda^{-1} \otimes x = f_{(\mathbf{K})}(\alpha^{-1}\lambda^{-1} \otimes x')\) , which proves the proposition.

COROLLARY 1. If \(\mathbf{E}'\) is a submodule of \(\mathbf{E}\) , \(\mathbf{E}_{(\mathbf{K})}'\) is canonically identified with a vector subspace of \(\mathbf{E}_{(\mathbf{K})}\) and \((\mathbf{E} / \mathbf{E}')_{(\mathbf{K})}\) with \(\mathbf{E}_{(\mathbf{K})} / \mathbf{E}_{(\mathbf{K})}'\) .

It suffices to apply Proposition 27 to the exact sequence

\[
0 \rightarrow \mathrm{E} ^ {\prime} \rightarrow \mathrm{E} \rightarrow \mathrm{E} / \mathrm{E} ^ {\prime} \rightarrow 0.
\]

COROLLARY 2. For every A-linear mapping \(f: \mathrm{E} \to \mathrm{F}\) , \(\operatorname{Ker}(f_{(\mathbf{K})}) = (\operatorname{Ker}(f))_{(\mathbf{K})}\) , \(\operatorname{Im}(f_{(\mathbf{K})}) = (\operatorname{Im}(f))_{(\mathbf{K})}\) , \(\operatorname{Coker}(f_{(\mathbf{K})}) = (\operatorname{Coker}(f))_{(\mathbf{K})}\) to within canonical isomorphisms. In particular, for \(f_{(\mathbf{K})}\) to be injective (resp. surjective, resp. zero), it is necessary and sufficient that \(\operatorname{Ker}(f) \subset \mathrm{T}(\mathrm{E})\) (resp. that \(\operatorname{Coker}(f)\) be a torsion module, resp. that \(\operatorname{Im}(f) \subset \mathrm{T}(\mathrm{F})\) ).

This follows from Corollary 1 and Proposition 26.

COROLLARY 3. Let E be an A-module and \((x_{\lambda})_{\lambda \in L}\) a family of elements of E. For \((x_{\lambda})\) to be a free family, it is necessary and sufficient that in the vector K-space \(\mathbf{E}_{(\mathbb{K})}\) the family \((1 \otimes x_{\lambda})\) be free.

The family \((x_{\lambda})\) defines an A-linear mapping \(f: \mathbf{A}^{(L)} \to \mathbf{E}\) such that \(f(e_{\lambda}) = x_{\lambda}\) for all \(\lambda \in \mathbf{L}((e_{\lambda})\) being the canonical basis of \(\mathbf{A}^{(L)})\) and to say that \((x_{\lambda})\) is free means that \(f\) is injective. It suffices to apply Corollary 2 to \(f\) , observing that \(\mathbf{A}^{(L)}\) is torsion-free (Proposition 25).

